\documentclass{book}
\usepackage{ctex}
\usepackage[cmex10]{amsmath}
\usepackage{xcolor}
\usepackage{mathtools}
\usepackage{booktabs}
\DeclareRobustCommand{\bluebold}[1]{\textbf{\textcolor[rgb]{0.31,0.63,0.86}{#1}}}

\usepackage[
    pdfpagemode=UseOutlines,
    bookmarks=true,
    bookmarksopen=true,
    bookmarksnumbered=true,
    colorlinks=true,
    linkcolor=black,
    citecolor=black,
    urlcolor=black
]{hyperref}
\usepackage{bookmark}

\raggedbottom

\begin{document}
\tableofcontents
\clearpage

化工

\begin{table}[htbp]
\centering
\begin{tabular}{c c c} 
\toprule
NO. & point & star\\ [0.5ex] 
\midrule
1 & HSE & 1\\
2 & 化学性质 & 1\\
\bottomrule
\end{tabular}
\caption{Knowledge points table for chemical.}
\label{table:pipeline-knowledge}
\end{table}

\chapter{HSE}

2026.11.13

丙烯，$\mathrm{H_2}$

HSE

\section{理念和目标}

\begin{itemize}
    \item 愿景：全球 HSE 领跑者（这就是自己建的，每个公司都有自己的这种企业文化）；
    \item 目标：零损失；
    \item 理念：HSE永远第一；
    \item 生命至上、环境优先；
    \item 损失控制、持续改进；
\end{itemize}

\section{要素}

\begin{itemize}
    \item 领导作用 规划保障 风险认知 能力建设 风险控制；
    \item 沟通推广 合规保证 风险监控 事件学习 持续改进；
\end{itemize}

要素1：领导作用

领导作用不可替代，优秀的领导力是HSE管理有效运行的核心驱动力。公司领导应做出明确的承诺，积极践行有感领导，发挥示范作用，建立实施FORUS体系所需的文化，用具体行动带领员工达成目标。

要素2：规划保障

规划指引正确的方向，资源是达成目标的必要保障。公司应在发展战略和行动计划中包含实现HSE愿景、目标的内容，为建立、实施和持续改进FORUS体系提供所需的资源。

要素3：风险认知

管理风险首先要认知风险，了解HSE风险及其后果对采取保护人员和环境所需的控制措施至关重要。公司应定期评估所有业务活动，以识别风险，确定合理有效的控制措施。

要素4：能力建设

HSE是被聘用的条件，员工应具备完成本职工作的基本能力。公司应建立良好的人力资源系统，对员工进行有效管理，并通过建立实施有效的培训体系，帮助员工了解工作场所及所从事活动的风险，始终有能力遵守正常运营和应对异常情况的管理程序。

要素5：风险控制

HSE管理的核心是控制风险。公司应根据风险控制的优先级别，采取合理可行的管控措施，使所有业务活动和过程的风险始终保持在可接受的水平。公司应在设计和建设阶段从本质上实现健康、安全、环保，在生产运营阶段不轻易改变原有意图，并通过有效的管理措施加以维持，如果出现异常或紧急情况则通过应急手段减少影响，使风险受控。

要素6：沟通推广

良好的沟通对确保所有员工理解、执行FORUS体系要求至关重要。公司应建立多样化的沟通渠道，确保重要信息有效的双向沟通，发现并广泛宣传突出的团队和个人表现以及最佳实践，强化正确的行为，同时与利益相关方保持互动，以获得认同。

要素7：合规保证

合法合规是中化香港集团HSE的底线。公司必须建立识别相关法律法规和标准规范的流程，评价其对业务产生的影响，以确保所有业务活动都能满足或高于法律法规和HSE管理要求，使主动合规成为公司的基本准则和自觉行为。

要素8：风险监控

监控风险才能识别改进机会。为了掌握风险是否受控，公司须建立风险监控机制，在认知风险与控制风险的基础上，利用HSE绩效数据监测和衡量HSE过程与趋势，识别风险控制不足，并采取必要的改进措施。

要素9：事件学习

有效的事件学习可以把不希望发生的事件转变成改进机会。公司应建设“无责备”文化，鼓励员工报告事件，开展事件调查，分析根本原因以采取有针对性的纠正和预防措施，定期进行事件统计分析和管理，以强化员工良好的HSE习惯。

要素10：持续改进

持续改进才能实现卓越。公司应通过总结分析、管理评审和管理对标，综合评估HSE绩效，发现良好实践和管理差距，并反馈到规划与计划制定中，持续提升管理水平。

\chapter{产品}

正丁醇、辛醇、异丁醇、异丁醛、新戊二醇

\end{document}
